% s-241-p241-m003; original PDF 241, 242
\Aphorism{s-241-p241-m003}{335}
Die Griechen, Staatsphilister und Politomanen wie sie waren, machen aus der Kunst eine Institution, eine Anstalt zu Staatszwecken, — dass wir das nicht nachahmen, scheint mir für \Sperr{unseren} besseren Geschmack zu sprechen. Die Künstler selber würden sich freilich für die griechische Kunstpflege erklären. Was redet aus allen zehn Bänden der Wagnerischen Kunstschriften, wenn nicht der Wille zur Institution! Bayreuth sollte eine Schule bedeuten (erst in zweiter Linie ein Theater), eine Stilfestlegung und Tradition für Jahrhunderte, ein aere perennius gegenüber den Tagesmoden — „Tag“ ist bei Wagner ein Schimpfwort, wie „Welt“ bei den Kirchenvätern. — Uns, den flüggen Wandervögeln und Betrachtern aller „Werke und Tage“, uns wird es nicht so leicht einfallen, die einmaligen Glückswürfe und Wunder der Kunst verlängert, verweilend, verlangweilt zu wünschen. Genug, wenn die Aloe wieder einmal geblüht hat! keine Treibhäuser für eine perennirende Staats- und Volkskunst! Es wird immer grosse Augenblicke und Gelegenheiten geben, und kein Abhang ist so steil und dürr, dass dort nicht irgend eine wunderliche Überraschung von Baum Wurzeln schlüge — aber uns will es nicht zu Sinn und Geschmack, ganze Alleen mit den Bäumen constanter Kunstproduction zu bepflanzen. \Greek{Δὶς} \Greek{καὶ} \Greek{τρὶς} \Greek{τὸ} \Greek{καλόν}: die Griechen, wie auch die Renaissance, zeigen am deutlichsten die Schattenseiten der Kunst als Institution: Wiederholung derselben Gegenstände bis zum Typischen, Erstarrung in Schule und Manier, Überschwemmung des Persönlichen durch die Fluth der Reminiscenzen und Vorbilder. Von dieser Gefahr kann sich keine unterstützte und garantirte Kunst fern halten; daneben vergesse man auch nicht, dass die Förderung durch den Staat — und wer käme sonst als dauernder Kunstpfleger in Betracht? — sehr bald in gewisse Forderungen von Seiten des Staates übergeht.
\par

